\begin{claim}{B.36}{Proved}{Complete exceptional catalogues and their realization fibers}
For the four exceptional types, the numbers of finite-subgroup
conjugacy classes in the extended mapping class group, and the
numbers of finite action classes inducing these subgroups
isomorphically, are as follows:
\[
\begin{array}{c|c|r|r}
(g,\nu)&\mathcal T_{g,\nu}&
 \text{finite subgroup classes}&\text{realized action classes}\\\hline
(0,0)&C_2&2&3\\
(0,1)&C_2&2&2\\
(0,2)&C_2^2&5&7\\
(1,0)&\mathrm{GL}_2(\Z)&13&17
\end{array}                                             \tag{B.36.1}
\]
The executed files
\path{exceptional_finite_subgroups_and_lifts.csv} and
\path{exceptional_sections.json} give all $29$ action rows over
these $22$ mapping class subgroup rows.  They retain the extension,
all full and partial section classes, $\Phi$ and its stabilizers,
centralizers, and every defined local multiplier.
The request in B(f) for one set of section data attached only to a
finite mapping class subgroup is false in these types: the displayed
realization fibers are required.  For the antipodal action on $S^2$,
the algebraic section exists and the fixed-point locus is empty;
there is no fixed-point rotation multiplier for that section.
\end{claim}
\begin{proof}
We first identify the groups in (B.36.1) and the comparison being
classified.  Orientation is the component invariant for
$\operatorname{Homeo}(S^2)$ and for
$\operatorname{Homeo}(S^2,\{\infty\})$.
An orientation-preserving homeomorphism is isotopic to the identity:
fix a disc and its complementary disc by the Alexander isotopy,
with the marked point kept fixed in the second case.
For the annulus, the two component invariants are the sign
$\epsilon$ on its fundamental group and the sign $\eta$ of its
permutation of the two ends.  Their four representatives are
\[
       (x,t)\longmapsto(\epsilon x,\eta t),
       \qquad x\in\R/\Z,\quad t\in\R.                  \tag{B.36.2}
\]
The kernel of these invariants is connected: compactify the two
ends, isotope an arc joining them to a fixed arc, and apply the
Alexander isotopy on the complementary disc.  Winding of the arc
is absorbed by rotating a puncture neighborhood, since its end
is a puncture and has no prescribed boundary parametrization.
For the torus, homology gives
$\pi_0\operatorname{Homeo}(T^2)=\mathrm{GL}_2(\Z)$.
Surjectivity is given by linear maps.  A homologically trivial
homeomorphism is first isotoped to fix the origin, then the two
basis curves, and then their complementary disc; translations
lie in the identity component.  B.25 gives its thirteen finite
subgroup classes.

In this statement a realized action over a finite subgroup
$H\subset\pi_0\operatorname{Homeo}(U)$ means a finite subgroup
$G\subset\operatorname{Homeo}(U)$ for which the projection
$G\to\pi_0\operatorname{Homeo}(U)$ is an isomorphism onto a conjugate
of $H$.  Actions are identified by homeomorphism conjugacy.
This specifies the exact fiber of the comparison map.  The
additional finite actions having a nontrivial kernel in that
projection are the full affine families already classified in
B.14; their kernels are retained in that result and are not
silently identified with an injective realization.

On $S^2$, a negative involution is conjugate either to a reflection
or to the antipodal map.  Indeed the sphere theorem in B.1 conjugates
the finite action into $O(3)$.  An orthogonal involution has
eigenvalues $\pm1$; negative determinant leaves precisely the
patterns $(1,1,-1)$ and $(-1,-1,-1)$.  Their fixed sets are a
circle and the empty set.  Together with the trivial action this
gives three rows over the two mapping class subgroups.
On the plane, a negative involution is conjugate to a reflection
by B.1, giving two rows including the trivial action.
For both surfaces $\pi=1$, so the lift extension is simply
$E=G$.  For the negative $C_2$ there is exactly one full section,
exactly one nontrivial partial section, and one maximal class,
with centralizer $1$ in $\pi$ and stabilizer $C_2$.
For the trivial group there is one full section and no nontrivial
partial section.  These are algebraic assertions about the
extension and include the antipodal action.  In that action
$p=-p$ implies $p=0$, which has no solution on the unit sphere.
The program verifies $A^2=I$, $\det A=-1$, and $\det(I-A)=8$
for $A=-I_3$.  It therefore records zero fixed components next
to the one algebraic full section.

A reflection has the local multiplier $1/2$.  For the antipodal
section no fixed-point germ exists from which to read such a
multiplier; the table marks that entry as undefined.  It separately
records the formal value $1/2$ obtained by extending the
orientation-sign convention to every negative involution.
This convention has a stated source and is not substituted for
a nonexistent fixed-point comparison.  In particular the
circle fixed by a sphere reflection has trivial image in
$\pi_1(S^2)$, so the hyperbolic oval-centralizer formula $\Z$
does not transfer to this case.

For the annulus use the simultaneous normal form of B.14.
Let
\[
 a(x,t)=(-x,t),\qquad
 b(x,t)=(x,-t),\qquad
 b'(x,t)=(x+\tfrac12,-t),\qquad
 c(x,t)=(-x,-t).                                      \tag{B.36.3}
\]
The three nontrivial mapping class elements are represented by
$a,b,c$.  A lift of the class of $a$ or $c$ has an arbitrary
angular shift, removed by translation conjugacy.
A lift of the class of $b$ has angular shift $0$ or $1/2$,
because its square is the identity.  These give $b$ and $b'$;
their fixed sets distinguish them.
For the whole four-group, normalize its angular reflection to
$a$.  Commutation forces the radial generator's angular shift
to be $0$ or $1/2$, giving exactly
$\langle a,b\rangle$ and $\langle a,b'\rangle$.
Conjugation cannot exchange the three nontrivial mapping class
elements, since their component group is abelian.  The fixed
oval or its absence again distinguishes these two actions.
This proves the seven-row classification.

Here are all its section counts; $f,p,\phi,o$ mean full sections,
nontrivial partial sections, $|\Phi|$, and $|\Phi/G|$:
\[
\begin{array}{l|rrrr}
\text{action}&f&p&\phi&o\\\hline
1&1&0&0&0\\
\langle a\rangle&2&2&2&2\\
\langle b\rangle&1&1&1&1\\
\langle b'\rangle&0&0&0&0\\
\langle c\rangle&2&2&2&2\\
\langle a,b\rangle&2&7&2&2\\
\langle a,b'\rangle&0&4&4&2
\end{array}                                             \tag{B.36.4}
\]
To verify them, the fixed sets of $a,b,b',c$ are respectively
two lines, one oval, the empty set, and two points.
In the split four-group, each of its two global fixed points
extends the corresponding point section and the incident
reflection sections; the oval section has both maximal extensions.
In the glide four-group, the two angular-reflection lines and the
two fixed points of $ab'$ are disjoint.  They are all maximal,
and the group interchanges each pair.  Thus (B.36.4) also gives
the required partial-section distinction.
Every nonidentity torsion element in these sections has multiplier
$1/2$.  The centralizer in $\pi=\Z$ is $\Z$ for the radial oval
section and $1$ for every point or line section.

These calculations can also be read directly from the extension.
For an affine element write its chosen angular shift as
$a_g\in[0,1)$ and let $\epsilon_g$ be its action on $\pi=\Z$.
The program records the integer cocycle
\[
 c(g,h)=a_g+\epsilon_ga_h-a_{gh}
 \quad\hbox{and}\quad
 (n,g)(m,h)=(n+\epsilon_gm+c(g,h),gh).                 \tag{B.36.5}
\]
It checks the cocycle identity on every triple.
In the universal cover $\R^2$, a finite subgroup of the lift
extension fixes a point; its fixed set is affine and connected.
Passing to the deck quotient identifies its conjugacy class
with the corresponding fixed component just listed.
This proves the section calculation without using a
center-free hypothesis.

For the torus apply the exact affine classification in B.27 and
retain precisely the normal forms whose translation kernel has
order one.  The linear projection of such a finite group is
faithful, which is exactly the comparison stipulated above.
Conversely every injective realized mapping class subgroup has
this property and is included by B.14.
There are seventeen affine forms over the thirteen integral
linear classes.  In the order of the linear classes in B.27,
the fiber sizes are
\[
                   1,2,1,1,1,3,1,1,1,1,1,2,1.       \tag{B.36.6}
\]
The proof of B.27 enumerates all cocycles and all affine
conjugacies without a word bound.  Its averaging comparison in
B.14 proves that no further topological conjugacy is missing.
Run the same exact fixed-point and reflection-circle computation
with no punctures removed.  Every isolated point section has
centralizer $1$ in $\Z^2$; every reflection oval has centralizer
the primitive invariant copy of $\Z$.  The trivial section has
centralizer $\Z^2$.  Local rotation matrices give the multipliers,
and the complete section poset gives the full and partial counts,
$\Phi$, and every stabilizer.  The resulting seventeen rows and
their exact rational matrices are included in the JSON.

The executed program \texttt{catalogue\_exceptional.py} carries
out all these finite calculations, records the two annulus-end
stabilizers and the plane's puncture stabilizer, and verifies the
four mapping class counts and the four realization-fiber totals.
Consequently it supplies a complete corrected computational
answer for the exceptional types, including the precise fibers
where B.22 prevents attaching one extension or one section table
to a bare mapping class subgroup.
\end{proof}
